\clearpage
\section{第 8.1 节：路径积分的规则}
\label{section:QM-C08-S01}
\studymeta
  {QM-C08-S01}
  {2026-09-23 至 2026-09-24}
  {Shankar, \textit{Principles of Quantum Mechanics}, Second Edition}
  {第 8 章 \S 8.1；正文印刷页 p.~223，PDF p.~234；图 8.1 位于印刷页 p.~224，PDF p.~235}
  {讲解与概念核对已完成；笔记已确认}

\textbf{范围与主线。}本节只给出路径积分（path integral）的基本规则：由连接固定时空端点的所有路径及其作用量，直接构造传播核。经典极限的分析留到第 8.2 节，时间切片的具体定义留到第 8.4 节。

\notetopic{QM-C08-S01-T01}{从作用量到传播核：对路径的振幅求和}
\textbf{来源：}教材式 (8.1.1) 与图 8.1。传播核的记号区分、概率解释及形式化求和的限定为学习补充。

\textbf{目标不变，求法改变。}用 $U(t_f,t_i)$ 表示抽象演化算符，用 $K$ 表示其位置表象的核；教材也以 $U$ 表示这个核：
\begin{align}
  K(x_f,t_f;x_i,t_i)&=\bra{x_f}U(t_f,t_i)\ket{x_i},\\
  \psi(x_f,t_f)&=\int_{-\infty}^{\infty}
     K(x_f,t_f;x_i,t_i)\psi(x_i,t_i)\dd x_i.
  \label{eq:QM-C08-S01-evolution}
\end{align}
原来的主要路线是从 Hamiltonian 的谱构造传播子；新路线从 Lagrangian 与作用量出发，直接求同一个传播核。求得 $K$ 后仍须按式~\eqref{eq:QM-C08-S01-evolution} 演化初态，再取 $|\psi(x_f,t_f)|^2$ 得到位置概率密度，不能把 $K$ 本身当作概率。

\textbf{固定端点，不限定中间路径满足运动方程。}以 $q(\tau)$ 表示一条候选路径，避免把路径函数与端点坐标混淆：
\begin{equation}
  q(t_i)=x_i,\qquad q(t_f)=x_f,\qquad
  S[q]=\int_{t_i}^{t_f}L(q(\tau),\dot q(\tau),\tau)\dd\tau.
\end{equation}
对于本节考虑的普通非相对论粒子，$L=\tfrac12m\dot q^{\,2}-V(q,\tau)$。不同路径具有相同的时空端点，但中间的位置、速度及作用量可以不同。求和并不要求每条路径满足 Newton 或 Euler--Lagrange 方程；经典路径只是其中的一条。

\begin{center}
\begin{minipage}{\linewidth}
\centering
\begin{tikzpicture}[>=Stealth,x=1cm,y=1cm]
  \draw[->] (0,0)--(8.3,0) node[right] {$\tau$};
  \draw[->] (0,0)--(0,2.65) node[above] {$q$};
  \coordinate (start) at (1,0.55);
  \coordinate (end) at (7,2.05);
  \draw[noteBlue,thick] (start) .. controls (2.6,2.3) and (5.2,2.5) .. (end);
  \draw[thick] (start) .. controls (2.8,0.6) and (5.5,1.4) .. (end);
  \draw[accentRed,thick] (start) .. controls (3.2,-0.15) and (5.4,0.4) .. (end);
  \draw[dashed,gray] (1,0)--(start);
  \draw[dashed,gray] (7,0)--(end);
  \fill (start) circle (2pt) node[left] {$x_i$};
  \fill (end) circle (2pt) node[right] {$x_f$};
  \node[below] at (1,0) {$t_i$};
  \node[below] at (7,0) {$t_f$};
\end{tikzpicture}

\small 固定两个时空端点，中间路径不同。依据教材图 8.1 的含义原创示意，曲线不指定具体势能。
\end{minipage}
\end{center}

\textbf{路径决定相位，振幅相加得到传播核。}教材给出的形式化规则为
\begin{equation}
  \boxed{K(x_f,t_f;x_i,t_i)
     =A\sum_{\substack{q(t_i)=x_i\\q(t_f)=x_f}}
         \ee^{\ii S[q]/\hbar}}.
  \label{eq:QM-C08-S01-recipe}
\end{equation}
求和遍及连接两端点的路径，不是只对经典解求和。作用量与 $\hbar$ 的量纲均为能量乘时间，因此 $S[q]/\hbar$ 是无量纲相位。对实作用量，$|\ee^{\ii S[q]/\hbar}|=1$；路径间的作用量差决定相位差，复振幅因而可能相增或相消。

\textbf{解释与适用边界。}
\begin{itemize}
  \item 相加的是\emph{概率振幅}，不是各路径的概率。相位因子模为 $1$ 不表示每条路径的发生概率为 $1$，也不能将规则解释为先随机选一条经典路径。
  \item $A$ 表示整体归一化因子；它须结合路径求和的测度与离散化理解，不能靠给每条路径分配普通概率来确定。
  \item 连续的路径集合不是普通的有限或可数求和。本节保留教材的形式化写法；时间切片将进一步说明其含义，不能据此假定所有路径都是光滑的经典轨迹。
\end{itemize}

\textbf{与已有知识的联系。}双缝实验要求将不可区分备选过程的振幅相加；这里把振幅叠加用于固定起终点之间的所有中间路径，并以作用量决定相位。这是理解规则的类比，不是从双缝实验或经典 $\delta S=0$ 单独推导出式~\eqref{eq:QM-C08-S01-recipe}。本节先陈述规则，后续再讨论经典极限及其与原有量子演化的联系。

\textbf{一分钟回顾。}固定时空端点，计算各路径的作用量，按 $\ee^{\ii S/\hbar}$ 叠加得到传播核；再让传播核作用于初始波函数。变化的是求传播子的途径，不是量子态与概率的解释。

\clearpage
\section{第 8.2 节：平稳相位与经典路径}
\label{section:QM-C08-S02}
\studymeta
  {QM-C08-S02}
  {2026-09-24 至 2026-09-25}
  {Shankar, \textit{Principles of Quantum Mechanics}, Second Edition}
  {第 8 章 \S 8.2；印刷页 pp.~224--225，PDF pp.~235--236；式 (8.2.1)--(8.2.2)，图 8.2--8.3}
  {讲解与教材正文例子已核对；笔记已确认}

\textbf{范围与主线。}各路径的相位因子模长相同，经典路径为何仍能在适当极限下占主导？本节用平稳相位解释相干叠加，再以自由粒子比较宏观与微观的作用量尺度。教材采用启发式路径求和，尚未证明它与 Schrödinger 演化等价；传播子的具体近似留到第 8.3 节。

\notetopic{QM-C08-S02-T01}{经典路径的优势来自邻近路径的相干叠加}
\textbf{来源：}教材印刷页 p.~224 及图 8.2；下列变分展开是联系第二章的推导补充。

对实作用量，每条路径贡献的相位因子满足
\begin{equation}
  Z[q]=\ee^{\ii S[q]/\hbar},\qquad |Z[q]|=1.
\end{equation}
经典路径并没有更大的单项模长。求和结果取决于各项的\emph{相位关系}：相位随路径快速变化时，复振幅大量相消；相位变化缓慢时，邻近的一簇路径可以较一致地相加。这就是 stationary phase（平稳相位）的基本图像；单项模长相同不等于每条路径具有相同的可观测概率。

\textbf{为什么经典路径附近相位更稳定？}令
\begin{equation}
  q_\epsilon(t)=q(t)+\epsilon\eta(t),\qquad
  \eta(t_i)=\eta(t_f)=0,
\end{equation}
其中 $\epsilon$ 是小的无量纲参数，$\eta$ 是保持端点不变的光滑变形。对足够光滑的作用量作局部展开：
\begin{equation}
  S[q+\epsilon\eta]
  =S[q]+\epsilon\,\delta S[q;\eta]
      +\frac{\epsilon^2}{2}\delta^2S[q;\eta]+O(\epsilon^3).
  \label{eq:QM-C08-S02-action-expansion}
\end{equation}
这里 $\delta S[q;\eta]=\left.\dd S[q+\epsilon\eta]/\dd\epsilon\right|_{\epsilon=0}$。由第二章的固定端点变分，
\begin{equation}
  \delta S[q;\eta]
  =\int_{t_i}^{t_f}\left(
      \frac{\partial L}{\partial q}
      -\frac{\dd}{\dd t}\frac{\partial L}{\partial\dot q}
    \right)\eta(t)\dd t.
\end{equation}
经典路径 $q_{\mathrm{cl}}$ 满足 Euler--Lagrange 方程，因而对\emph{所有}允许的 $\eta$，一阶变分均为零。于是
\begin{equation}
  \Delta\theta
  =\frac{S[q_{\mathrm{cl}}+\epsilon\eta]-S[q_{\mathrm{cl}}]}{\hbar}
  =\frac{\epsilon^2}{2\hbar}\delta^2S[q_{\mathrm{cl}};\eta]
       +O(\epsilon^3).
\end{equation}
相位变化从二阶或更高阶开始，而非一般的一阶。\emph{平稳不等于完全不变，也不要求作用量一定取极小值。}能较一致地贡献的是经典路径附近的一簇路径，不是只剩一条经典路径。

\begin{center}
\begin{minipage}{\linewidth}
\centering
\begin{tikzpicture}[>=Stealth,x=0.85cm,y=0.85cm]
  \begin{scope}
    \draw[->,gray] (0,0)--(4.9,0) node[right] {$\operatorname{Re}Z$};
    \draw[->,gray] (0,0)--(0,1.7) node[above] {$\operatorname{Im}Z$};
    \draw[->,thick,noteBlue] (0,0)--++(12:1) coordinate (a1);
    \draw[->,thick,noteBlue] (a1)--++(15:1) coordinate (a2);
    \draw[->,thick,noteBlue] (a2)--++(18:1) coordinate (a3);
    \draw[->,thick,noteBlue] (a3)--++(21:1);
    \node at (2,-0.6) {相位接近：合振幅较大};
  \end{scope}
  \begin{scope}[xshift=7cm]
    \draw[->,gray] (0,0)--(3.6,0) node[right] {$\operatorname{Re}Z$};
    \draw[->,gray] (0,0)--(0,1.7) node[above] {$\operatorname{Im}Z$};
    \draw[->,thick,accentRed] (0,0)--++(20:1) coordinate (b1);
    \draw[->,thick,accentRed] (b1)--++(105:1) coordinate (b2);
    \draw[->,thick,accentRed] (b2)--++(190:1) coordinate (b3);
    \draw[->,thick,accentRed] (b3)--++(275:1);
    \node at (1.5,-0.6) {相位分散：大量抵消};
  \end{scope}
\end{tikzpicture}

\small 依据教材图 8.2 的含义原创示意：等长复振幅首尾相接。箭头方向仅示意相位关系，并非实际计算的路径积分。
\end{minipage}
\end{center}

\notetopic{QM-C08-S02-T02}{作用量差与相干范围}
\textbf{来源：}教材印刷页 pp.~224--225 的相干范围估计与自由粒子例子。

两条路径的相对相位由作用量\emph{差}决定：
\begin{equation}
  \Delta\theta=\frac{\Delta S}{\hbar},\qquad
  \Delta S=S[q]-S[q_{\mathrm{cl}}].
\end{equation}
作用量与 $\hbar$ 的量纲均为能量乘时间，故该比值无量纲。教材用
\begin{equation}
  \frac{|\Delta S|}{\hbar}\lesssim\pi
\end{equation}
粗略估计经典路径附近的相干范围。它不是硬边界，也不是判定某条路径是否存在的规则。在二阶项主导时，二阶变分的绝对值越大，在给定变形方向上保持较小相位变化所允许的偏离就越小。

\textbf{经典极限的含义。}在相关作用量变化尺度远大于 $\hbar$、且平稳相位近似适用时，非平稳区域的振幅随路径快速振荡，求和后大量相消；经典路径附近的贡献因而突出。若有多条经典路径，可能需要保留多个平稳区域，不能预先假定只有一个。

不能仅凭任意两条路径的 $|\Delta S|/\hbar$ 很大就断言它们相消，因为相位以 $2\pi$ 为周期。关键是\emph{相位在一片路径中变化得快不快，以及整体求和的结果}。因此经典极限不是逐条删除非经典路径，也不能仅凭质量大小判断。教材中质量、距离和传播时间共同决定相位变化的具体计算见\relatedexercise{QM-C08-S02-E01}{两条自由粒子路径的作用量比较}。

\textbf{一分钟回顾。}经典路径没有更大的单项振幅；它的特殊性在于一阶作用量变化为零，使附近一簇路径较易相干叠加。经典近似取决于作用量差相对于 $\hbar$ 的尺度及整体的振荡相消，而非相位完全不变或非经典路径不再参与。
